\chapter{花生古诗词-1}

\section{野望}

\subsection*{野望}
\textbf{【唐】王绩}

东皋薄暮望，徙倚欲何依。\\
树树皆秋色，山山唯落晖。\\
牧人驱犊返，猎马带禽归。\\
相顾无相识，长歌怀采薇。

\bigskip
\noindent\textbf{译文：}\\
黄昏的时候伫立在房舍东边的高地怅望，徘徊不定不知该归依何方。\\
每棵树都染上秋天的色彩，重重山岭都披覆着落日的余光。\\
放牧的人驱赶着牛群回家，猎人骑着马带着猎物各自返回。\\
大家相对无言彼此互不相识，只能咏一曲长歌来怀念古代采薇而食的隐士。

% \rule{\textwidth}{0.4pt}
% \bigskip

\subsection*{黄鹤楼}
\textbf{【唐】崔颢}

昔人已乘黄鹤去，此地空余黄鹤楼。\\
黄鹤一去不复返，白云千载空悠悠。\\
晴川历历汉阳树，芳草萋萋鹦鹉洲。\\
日暮乡关何处是？烟波江上使人愁。

\bigskip
\noindent\textbf{译文：}\\
过去的仙人已经驾着黄鹤飞走了，只留下空荡荡的黄鹤楼。\\
黄鹤一去再也没有回来，千百年来只看见白云在天上飘浮。\\
阳光照耀下的汉阳树木清晰可见，更能看清芳草萋萋的鹦鹉洲。\\
暮色渐渐漫起，哪里是我的家乡？江面烟波渺渺让人更生忧愁。

% \rule{\textwidth}{0.4pt}
% \bigskip

\subsection*{使至塞上}
\textbf{【唐】王维}

单车欲问边，属国过居延。\\
征蓬出汉塞，归雁入胡天。\\
大漠孤烟直，长河落日圆。\\
萧关逢候骑，都护在燕然。

\bigskip
\noindent\textbf{译文：}\\
轻车简从将要去慰问边关，我要到远在西北边塞的居延。\\
像随风而去的蓬草一样出临边塞，北归大雁正翱翔云天。\\
浩瀚沙漠中孤烟直上云霄，黄河边上落日正圆。\\
到萧关时遇到侦察骑兵，得知主帅尚在前线未归。


\subsection*{渡荆门送别}
\textbf{【唐】李白}

渡远荆门外，来从楚国游。\\
山随平野尽，江入大荒流。\\
月下飞天镜，云生结海楼。\\
仍怜故乡水，万里送行舟。

\bigskip
\noindent\textbf{译文：}\\
我乘舟渡过长江来到遥远的荆门外，来到战国时期楚国的境内游览。\\
山随着平坦广阔的原野的出现逐渐消失，江水在一望无际的原野中奔流。\\
江面月影好似天上飞来的明镜，云层倒映构成海市蜃楼。\\
我依然喜爱这来自故乡之水，不远万里送我东行的舟儿。

% \rule{\textwidth}{0.4pt}
% \bigskip

\subsection*{钱塘湖春行}
\textbf{【唐】白居易}

孤山寺北贾亭西，水面初平云脚低。\\
几处早莺争暖树，谁家新燕啄春泥。\\
乱花渐欲迷人眼，浅草才能没马蹄。\\
最爱湖东行不足，绿杨阴里白沙堤。

\bigskip
\noindent\textbf{译文：}\\
从孤山寺的北面到贾亭的西面，湖面春水刚与堤平，白云低垂，同湖面上连成一片。\\
几只早出的黄莺争相飞往向阳的树木，谁家新飞来的燕子忙着筑巢衔泥。\\
纷繁的花朵渐渐开放使人眼花缭乱，浅浅的青草刚刚够上遮没马蹄。\\
最爱的湖东美景百游不厌，杨柳成排绿荫中穿过一条白沙堤。

% \rule{\textwidth}{0.4pt}
% \bigskip

\subsection*{庭中有奇树}
\textbf{【两汉】佚名}

庭中有奇树，绿叶发华滋。\\
攀条折其荣，将以遗所思。\\
馨香盈怀袖，路远莫致之。\\
此物何足贵？但感别经时。

\bigskip
\noindent\textbf{译文：}\\
庭院里一株佳美的树，满树绿叶托着繁盛的花朵。\\
我攀着树枝，摘下了其中一朵，想把它赠送给心中日夜思念的人。\\
花香充满了我的衣服襟袖之间，可是天遥地远，没能送到心上人的手中。\\
并不是此花有什么珍贵，只是有感于离别多时，想借着花儿表达思念之情罢了。


\subsection*{赠从弟·其二}
\textbf{【魏晋】刘桢}

亭亭山上松，瑟瑟谷中风。\\
风声一何盛，松枝一何劲。\\
冰霜正惨凄，终岁常端正。\\
岂不罹凝寒？松柏有本性。

\bigskip
\noindent\textbf{译文：}\\
高山上松树挺拔耸立，山谷间狂风瑟瑟呼啸。\\
风声是多么的猛烈，松枝又是多么的刚劲！\\
任它满天冰霜惨惨凄凄，松树的腰杆终年端端正正。\\
难道是松树没有遭到严寒的侵凌吗？不，是松柏天生有着耐寒的本性！

% \rule{\textwidth}{0.4pt}
% \bigskip

\subsection*{梁甫行}
\textbf{【魏晋】曹植}

八方各异气，千里殊风雨。\\
剧哉边海民，寄身于草野。\\
妻子象禽兽，行止依林阻。\\
柴门何萧条，狐兔翔我宇。

\bigskip
\noindent\textbf{译文：}\\
八方的气候各不相同，千里之内的风雨形态不一。\\
海边的贫民多么艰苦啊，平时就住在野外的草棚里。\\
妻子和儿女像禽兽一样生活，盘桓在险阻的山林里。\\
简陋的柴门如此冷清，狐兔在房屋周围自在地行走毫无顾忌。

% \rule{\textwidth}{0.4pt}
% \bigskip

\subsection*{龟虽寿}
\textbf{【两汉】曹操}

神龟虽寿，犹有竟时。\\
腾蛇乘雾，终为土灰。\\
老骥伏枥，志在千里。\\
烈士暮年，壮心不已。\\
盈缩之期，不但在天；\\
养怡之福，可得永年。\\
幸甚至哉，歌以咏志。

\bigskip
\noindent\textbf{译文：}\\
神龟虽能长寿，但也有死亡的时候。\\
腾蛇尽管能乘雾飞行，终究也会死亡化为土灰。\\
年老的千里马虽然伏在马槽旁，它的雄心壮志仍然是能够驰骋千里。\\
有远大抱负的人士到了晚年，奋发思进的雄心不会止息。\\
人的寿命长短，不只是由上天所决定的。\\
只要自己调养好身心，也可以益寿延年。\\
啊，庆幸得很！就用诗歌来表达内心的志向吧！



\subsection*{饮酒·其五}
\textbf{【魏晋】陶渊明}

结庐在人境，而无车马喧。\\
问君何能尔？心远地自偏。\\
采菊东篱下，悠然见南山。\\
山气日夕佳，飞鸟相与还。\\
此中有真意，欲辨已忘言。

\bigskip
\noindent\textbf{译文：}\\
将房屋建造在人来人往的地方，却不会受到世俗交往的喧扰。\\
问我为什么能这样，只要心中所想远离世俗，自然就会觉得所处地方僻静了。\\
在东篱之下采摘菊花，悠然间，那远处的南山映入眼帘。\\
傍晚时分南山景致甚佳，雾气峰间缭绕，飞鸟结伴而还。\\
这里面蕴含着人生的真正意义，想要分辨清楚，却已忘了怎样表达。

% \rule{\textwidth}{0.4pt}
% \bigskip

\subsection*{春望}
\textbf{【唐】杜甫}

国破山河在，城春草木深。\\
感时花溅泪，恨别鸟惊心。\\
烽火连三月，家书抵万金。\\
白头搔更短，浑欲不胜簪。

\bigskip
\noindent\textbf{译文：}\\
国都遭破旧山河依旧，长安城里的杂草和树木茂盛地荒芜。\\
感伤国事，看到花开不禁潸然泪下，内心惆怅怨恨，听到鸟鸣，竟觉得那叫声也很揪心。\\
连绵的战火已经延续了一个春天，家书难得，一封抵得上万两黄金。\\
愁绪缠绕，搔头思考，白发越搔越短，简直要不能插簪了。

% \rule{\textwidth}{0.4pt}
% \bigskip

\subsection*{雁门太守行}
\textbf{【唐】李贺}

黑云压城城欲摧，甲光向日金鳞开。\\
角声满天秋色里，塞上燕脂凝夜紫。\\
半卷红旗临易水，霜重鼓寒声不起。\\
报君黄金台上意，提携玉龙为君死。

\bigskip
\noindent\textbf{译文：}\\
敌兵滚滚而来，犹如黑云翻卷，想要摧倒城墙；战士们的铠甲在阳光照射下金光闪烁。\\
号角声响彻秋夜的长空，边塞上将士的血迹在寒夜中凝为紫色。\\
红旗半卷，援军赶赴易水；夜寒霜重，鼓声悲壮。\\
为了报答国君的赏赐和厚爱，手挥宝剑甘愿为国血战到底。


\subsection*{渔家傲}
\textbf{【宋】李清照}

天接云涛连晓雾，星河欲转千帆舞。\\
仿佛梦魂归帝所。闻天语，殷勤问我归何处。\\
我报路长嗟日暮，学诗漫有惊人句。\\
九万里风鹏正举。风休住，蓬舟吹取三山去！

\bigskip
\noindent\textbf{译文：}\\
水天相接，晨雾蒙蒙笼云涛。银河转动，像无数的船只在舞动风帆。\\
梦魂仿佛回到天庭，听见天帝在对我说话。他热情而又充满诚意地问我到哪里去。\\
我回答天帝路途还很漫长，现在已是黄昏却还未到达。即使我学诗能写出惊人的句子，又有什么用呢？\\
长空九万里，大鹏冲天飞正高。风啊！千万别停息，将我这一叶轻舟，直送往蓬莱三仙岛。

% \rule{\textwidth}{0.4pt}
% \bigskip

\subsection*{月夜}
\textbf{【唐】杜甫}

今夜鄜州月，闺中只独看。\\
遥怜小儿女，未解忆长安。\\
香雾云鬟湿，清辉玉臂寒。\\
何时倚虚幌，双照泪痕干。

\bigskip
\noindent\textbf{译文：}\\
今夜鄜州上空那轮圆月，只有你在闺房中独自遥看。\\
远在他乡怜惜幼小的儿女，还不懂得你为何思念长安？\\
蒙蒙雾气沾湿了你的鬓发；清冷的月光使你的玉臂生寒。\\
什么时候才能在一起共同靠在透光的窗帘旁，让月光擦干两人的思念的泪。

% \rule{\textwidth}{0.4pt}
% \bigskip

\subsection*{赤壁}
\textbf{【唐】杜牧}

折戟沉沙铁未销，自将磨洗认前朝。\\
东风不与周郎便，铜雀春深锁二乔。

\bigskip
\noindent\textbf{译文：}\\
赤壁的泥沙中，埋着一枚未锈尽的断戟。自己磨洗后发现这是当年赤壁之战的遗物。\\
倘若不是东风给周瑜以方便，结局恐怕是曹操取胜，二乔被关进铜雀台了。

% \rule{\textwidth}{0.4pt}
% \bigskip

\subsection*{浣溪沙}
\textbf{【宋】晏殊}

一曲新词酒一杯，去年天气旧亭台。夕阳西下几时回？\\
无可奈何花落去，似曾相识燕归来。小园香径独徘徊。

\bigskip
\noindent\textbf{译文：}\\
填一曲新词喝一杯美酒，还是去年的天气旧日的亭台，西落的夕阳何时才能回来？\\
花儿总要凋落让人无可奈何，似曾相识的春燕又归来，独自在花香小径里徘徊留恋。

\end{document}